%% AI4SE Project 1 -- Issue Report Classification (NLBSE'24)
%% Compile with:  make report        (latexmk -pdf report.tex)
%%
%% Numbers are NEVER typed by hand. Every table in section 5 is generated by
%%     python experiments/04_make_tables.py
%% which reads results/*.json and writes report/tables/*.tex. If a table is
%% missing, the experiment that produces it has not been run yet -- the
%% placeholder says which one.

\documentclass[10pt,a4paper]{article}

\usepackage[utf8]{inputenc}
\usepackage[T1]{fontenc}
\usepackage[margin=2.4cm]{geometry}
\usepackage{booktabs}
\usepackage{graphicx}
\usepackage{amsmath}
\usepackage{siunitx}
\usepackage{xcolor}
\usepackage{listings}
\usepackage[hidelinks]{hyperref}
\usepackage{caption}

\sisetup{round-mode=places, round-precision=4, detect-weight=true}

% Marks a section the team still has to write. Every one of these must be gone
% before submission; searching the source for "TODO" is the final check.
\newcommand{\todoitem}[1]{\textcolor{red}{\textbf{[TODO]} #1}}

% Included when an experiment has not produced its table yet.
\newcommand{\missingtable}[1]{%
  \begin{center}\ttfamily\small
  [table not generated --- run \detokenize{#1}]
  \end{center}}

\lstset{
  basicstyle=\ttfamily\footnotesize,
  breaklines=true,
  frame=single,
  showstringspaces=false,
}

\title{Issue Report Classification on the NLBSE'24 Benchmark:\\
       Baselines, Reproduction, and the Limits of the Evaluation}
\author{%
  \todoitem{names and matriculation numbers} \\
  \small Artificial Intelligence for Software Engineering ---
  University of L'Aquila
}
\date{\today}

\begin{document}
\maketitle

\begin{abstract}
We address the NLBSE'24 tool competition task: classifying GitHub issue
reports as \texttt{bug}, \texttt{feature} or \texttt{question} across five
open-source projects, with one classifier trained per project. We implement a
persistence-agnostic data layer, a range of classical and neural classifiers,
and a reproduction of the organisers' SetFit baseline.

Beyond the leaderboard, we report two properties of the benchmark itself.
First, its test set cannot resolve differences below approximately three F1
points, while re-running the same configuration with a different random seed
already moves the score by up to 2.4 points --- so a large share of the
differences reported on this task are not measurable. Second, class labels and
issue creation dates are confounded: a classifier given only the timestamp,
and no text at all, reaches an F1 of 0.68 against a chance level of 0.33, and
evaluating chronologically rather than randomly costs roughly eight points. We
argue that results on this benchmark should be accompanied by a power analysis
and a time-aware evaluation, and we release the code to do both.
\end{abstract}

%======================================================================
\section{Introduction}
\label{sec:intro}

\todoitem{Two or three paragraphs: why automatic issue triage matters, what
the NLBSE'24 competition asks for, and what this report contributes. Keep the
contributions list below in sync with what actually got done.}

Our contributions are:
\begin{enumerate}
  \item A modular implementation in which the persistence layer (in-memory or
        file-backed) can be exchanged without altering any business logic,
        verified by a test that runs the same pipeline over both.
  \item Classical and neural baselines, each selected by $k$-fold
        cross-validation on the training split and evaluated once on the test
        split over five random seeds.
  \item A faithful reproduction of the organisers' SetFit baseline.
  \item A power analysis of the benchmark, establishing the smallest
        difference its test set can detect.
  \item A time-aware evaluation protocol that quantifies the optimism of the
        official random split.
\end{enumerate}

%======================================================================
\section{Related work}
\label{sec:related}

\todoitem{One page. Anchors below; expand each into two or three sentences and
add the NLBSE'23 and '24 competition reports.}

Automatic issue classification was framed by Antoniol et al.\ and revisited at
scale by Kallis et al., whose Ticket Tagger applies fastText to the problem and
is the direct ancestor of this competition~\cite{kallis2019tickettagger}.
Colavito et al.\ introduced the SetFit approach that the organisers adopted as
the baseline, and reported a notable gap between training on the raw labels and
on a manually verified subset~\cite{colavito2023fewshot}. That gap is
consistent with Herzig et al., who manually inspected more than seven thousand
issue reports and found 33.8\% of those labelled as bugs to be
misclassified~\cite{herzig2013misclassification}.

Our methodological arguments follow the software-engineering literature on
evaluation validity: the need to preserve temporal order when validating
models trained on evolving repositories, and the use of Bayesian analysis for
comparing classifiers where frequentist tests cannot express
equivalence~\cite{benavoli2017bayesian}.

%======================================================================
\section{Dataset}
\label{sec:dataset}

The competition provides 3{,}000 labelled issues drawn from
\texttt{facebook/react}, \texttt{tensorflow/tensorflow},
\texttt{microsoft/vscode}, \texttt{bitcoin/bitcoin} and
\texttt{opencv/opencv}, split 50/50 into training and test sets and balanced at
exactly 100 issues per (project, class) cell in each split. Issues carrying
more than one label were removed by the organisers, so the task is multi-class
with a single label per instance, not multi-label.

\begin{table}[htbp]
\centering
\caption{Composition of each split. Both splits have identical structure.}
\label{tab:dataset}
\begin{tabular}{lrrrrl}
\toprule
Project & bug & feature & question & total & Issue dates \\
\midrule
\texttt{facebook/react} & 100 & 100 & 100 & 300 & 2016-03 -- 2023-08 \\
\texttt{tensorflow/tensorflow} & 100 & 100 & 100 & 300 & 2021-12 -- 2023-09 \\
\texttt{microsoft/vscode} & 100 & 100 & 100 & 300 & 2023-06 -- 2023-09 \\
\texttt{bitcoin/bitcoin} & 100 & 100 & 100 & 300 & 2018-10 -- 2023-09 \\
\texttt{opencv/opencv} & 100 & 100 & 100 & 300 & 2022-01 -- 2023-09 \\
\midrule
total & 500 & 500 & 500 & 1500 & \\
\bottomrule
\end{tabular}
\caption*{\footnotesize Date ranges measured from the CSVs and span both splits. Three of the five projects begin well before the January 2022 start date given in the competition materials.}
\end{table}


\subsection{Characterisation}

Median length is about 150 whitespace tokens, with a long tail: the longest
issue exceeds 21{,}000. At a 512-token cut-off, 97.3\% of \texttt{feature}
issues, 93.4\% of \texttt{question} and 91.9\% of \texttt{bug} issues are
retained in full, which is what motivates our choice of truncation length in
Section~\ref{sec:models}.

Issue bodies are Markdown documents, and their structure is not noise. A
logistic regression given only fifteen boolean structural flags --- does the
body contain a fenced code block, a stack trace, an image, a leftover HTML
comment from the issue template, a question mark in the title --- reaches an F1
of 0.556 without reading a single word. Stack traces appear in 31.9\% of
\texttt{bug} issues against 8.0\% of \texttt{feature} ones; unremoved template
comments appear in 29.5\% of \texttt{question} issues against 9.5\% of
\texttt{bug} ones. A permutation control confirms this is signal rather than
model capacity: with shuffled labels the same model scores
$0.340 \pm 0.039$ over twenty permutations ($z = 5.5$).

\subsection{A temporal confound}
\label{sec:confound}

In all five projects, the median creation date of \texttt{bug} issues is later
than that of the other two classes; the gap between the earliest and latest
class median ranges from 50 days in \texttt{microsoft/vscode} to 1{,}522 days
in \texttt{facebook/react}. This follows from how the dataset was built:
sampling a fixed 100 issues per class from repositories whose class prevalence
drifts over time makes the sampling date informative about the class.

\begin{table}[htbp]
\centering
\caption{F1 of a classifier given only \texttt{created\_at} and no text.}
\label{tab:timestamp}
\begin{tabular}{lr}
\toprule
Project & F1 from timestamp alone \\
\midrule
\texttt{facebook/react} & 0.802 \\
\texttt{tensorflow/tensorflow} & 0.792 \\
\texttt{bitcoin/bitcoin} & 0.622 \\
\texttt{opencv/opencv} & 0.620 \\
\texttt{microsoft/vscode} & 0.587 \\
\midrule
mean & 0.684 \\
\bottomrule
\end{tabular}
\caption*{\footnotesize Chance level is 0.333. Five-fold cross-validation within each project.}
\end{table}


The effect is not removable by dropping the timestamp column. Adding
\texttt{created\_at} as an explicit feature to a text model changes its score
by 0.0006, because the information is already carried by era-specific
vocabulary: the Jensen--Shannon divergence between the vocabulary of the older
and newer halves of the data is 0.1475.

%======================================================================
\section{Approach}
\label{sec:approach}

\subsection{Architecture and persistence}
\label{sec:architecture}

The project specification requires that data be managed both in memory and
through files, and that switching between the two involve few changes to the
application. We implement the Repository pattern: an abstract
\texttt{IssueRepository} declares five primitives (\texttt{load}, \texttt{save},
\texttt{all}, \texttt{add}, \texttt{clear}), and every query helper used by the
rest of the system is implemented once on the abstract base class in terms of
those five. Two concrete subclasses back the collection with a Python list and
with a CSV or JSON file respectively.

Switching layer is a single argument:

\begin{lstlisting}[language=Python]
train = load_split("train", kind="memory")  # held in RAM
train = load_split("train", kind="file")    # backed by a CSV on disk
\end{lstlisting}

Nothing downstream changes, because nothing downstream refers to a concrete
repository class. \texttt{tests/test\_persistence\_equivalence.py} runs the same
analysis pipeline over both layers and asserts the outputs are identical.

\todoitem{Add the class diagram. \texttt{results/figures/} is the agreed
location.}

\subsection{Preprocessing}

Three cleaning levels are provided so that the effect of cleaning can be
ablated by changing one string: \texttt{raw} (whitespace only), \texttt{light}
(removes code, stack traces, markup, URLs, paths, SHAs and mentions) and
\texttt{full} (\texttt{light} plus lowercasing, punctuation removal, stop-word
removal and lemmatisation). Average retained length falls to 63\% and 33\% of
the original respectively.

Given the structural findings of Section~\ref{sec:dataset}, deletion is not
obviously the right choice, and we treat the cleaning level as a
hyper-parameter rather than a fixed preprocessing decision.

\subsection{Classifiers}
\label{sec:models}

\todoitem{One short paragraph per family. The implementations live in
\texttt{src/ai4se/classifiers/}.}

\begin{description}
  \item[Classical] TF-IDF word $n$-grams feeding multinomial Naive Bayes,
        logistic regression, a linear SVM and a random forest.
  \item[Neural] A two-layer perceptron over TF-IDF features, and the
        convolutional sentence classifier of Kim with embeddings learned from
        scratch.
  \item[SetFit] The organisers' baseline: \texttt{all-mpnet-base-v2}
        fine-tuned contrastively, then a logistic-regression head.
  \item[Transformer] \todoitem{Fine-tuned encoder, if run. Report honestly if
        it was not.}
\end{description}

%======================================================================
\section{Evaluation protocol}
\label{sec:protocol}

\subsection{Metric}

The competition ranks submissions by the mean of five per-project F1 scores,
where each project's score is the support-weighted average over the three
classes. Because every cell contains exactly 100 issues, that weighted average
coincides with the unweighted one. We report this figure throughout, alongside
the full per-project, per-class table --- the single number hides a range from
0.62 to 0.92 across those cells.

\subsection{Model selection}

All hyper-parameters are selected by 5-fold cross-validation stratified on
(project, class) pairs, \emph{using the training split only}. Within each fold
one model is trained per project, mirroring the final protocol rather than a
simplification of it. The test split is read exactly once per configuration,
at the end.

We deliberately do not use a project-grouped split for selection: grouping by
project removes all in-domain data and turns every fold into a
leave-one-project-out problem, which scores 0.617 and is not the quantity the
competition measures.

\subsection{Uncertainty, significance and power}
\label{sec:power}

Each configuration is trained with five seeds and reported as mean $\pm$
standard deviation, with all five raw values given in
Appendix~\ref{sec:appendix-seeds}. Confidence intervals come from a bootstrap
stratified by project, $B = 10{,}000$. Paired comparisons use the exact
binomial form of McNemar's test, computed per project and corrected across the
family of five with the Holm--Bonferroni procedure.

The reason for that machinery is quantified in
Table~\ref{tab:power}: simulation of McNemar's test at the observed discordant
rate gives a minimum detectable difference of approximately three F1 points on
the full test set, rising to seven points within a single project. Seed
variation alone spans 2.4 points. We therefore adopt the following rule, fixed
in \texttt{preregistration.md} before any model was trained: \emph{a difference
below three points is reported as ``no detectable difference'' and is never
described as an improvement.}

\begin{table}[htbp]
\centering
\caption{Minimum detectable difference at 80\% power, $\alpha = 0.05$.}
\label{tab:power}
\begin{tabular}{lr}
\toprule
Evaluation unit & Detectable difference (F1 points) \\
\midrule
Full test set (1500 issues) & 3.0 \\
One project (300 issues) & 7.0 \\
\midrule
Seed variation alone, same configuration & 2.4 \\
\bottomrule
\end{tabular}
\caption*{\footnotesize Monte Carlo simulation of McNemar's test at the observed discordant rate of 0.161.}
\end{table}


%======================================================================
\section{Results}
\label{sec:results}

\subsection{Baseline reproduction}

\missingtable{experiments/03_setfit.py}


\subsection{Main comparison}

\begin{table}[htbp]
\centering
\caption{Cross-repository F1. Our models are the mean of five seeds.}
\label{tab:main}
\begin{tabular}{lrrc}
\toprule
Model & F1 & sd (seeds) & 95\% CI \\
\midrule
setfit (published) & 0.8240 & --- & --- \\
roberta (published) & 0.7923 & --- & --- \\
fasttext (published) & 0.7184 & --- & --- \\
\midrule
linear\_svm & 0.7504 & 0.0000 & [0.729, 0.773] \\
logreg & 0.7502 & 0.0000 & [0.729, 0.772] \\
ffnn & 0.7362 & 0.0043 & [0.714, 0.758] \\
text\_cnn & 0.7255 & 0.0067 & [0.703, 0.748] \\
random\_forest & 0.7241 & 0.0010 & [0.702, 0.746] \\
naive\_bayes & 0.7199 & 0.0000 & [0.696, 0.743] \\
\bottomrule
\end{tabular}
\caption*{\footnotesize Intervals are stratified bootstraps on the median-seed run. Differences below 3.0 points are not detectable on this test set (Table~\ref{tab:power}) and are not described as improvements.}
\end{table}


\subsection{Per-project and per-class detail}

\begin{table}[htbp]
\centering
\caption{Per-project, per-class F1 of the strongest model (\texttt{linear\_svm}).}
\label{tab:perrepoclass}
\begin{tabular}{lrrr}
\toprule
Project & bug & feature & question \\
\midrule
\texttt{facebook/react} & 0.925 & 0.765 & 0.674 \\
\texttt{tensorflow/tensorflow} & 0.817 & 0.811 & 0.703 \\
\texttt{microsoft/vscode} & 0.737 & 0.696 & 0.739 \\
\texttt{bitcoin/bitcoin} & 0.667 & 0.815 & 0.657 \\
\texttt{opencv/opencv} & 0.670 & 0.788 & 0.794 \\
\midrule
mean & 0.763 & 0.775 & 0.713 \\
\bottomrule
\end{tabular}
\caption*{\footnotesize The single cross-repository figure conceals this spread; the weakest and strongest cells differ by roughly 30 points.}
\end{table}


\subsection{Significance}

\begin{table}[htbp]
\centering
\caption{Per-project model against pooled model, exact McNemar.}
\label{tab:significance}
\begin{tabular}{lrrrrrc}
\toprule
Project & $n_{01}$ & $n_{10}$ & discordant & $p$ & $p_{\text{Holm}}$ & sig. \\
\midrule
\texttt{bitcoin/bitcoin} & 17 & 9 & 26 & 0.169 & 0.843 & ns \\
\texttt{facebook/react} & 14 & 17 & 31 & 0.720 & 1.000 & ns \\
\texttt{microsoft/vscode} & 17 & 10 & 27 & 0.248 & 0.991 & ns \\
\texttt{opencv/opencv} & 16 & 14 & 30 & 0.856 & 1.000 & ns \\
\texttt{tensorflow/tensorflow} & 21 & 15 & 36 & 0.405 & 1.000 & ns \\
pooled & 85 & 65 & 150 & 0.121 & 0.121 & ns \\
\bottomrule
\end{tabular}
\caption*{\footnotesize 0 of five projects reach significance after correction. The two configurations are not distinguishable on this test set.}
\end{table}


\subsection{Time-aware evaluation}
\label{sec:timeaware}

\begin{table}[htbp]
\centering
\caption{Cost of evaluating on a random rather than a chronological split.}
\label{tab:timeaware}
\begin{tabular}{lr}
\toprule
Protocol & Cross-repository F1 \\
\midrule
Random split (matched control) & 0.7700 \\
Time-aware split & 0.6817 \\
\midrule
Difference & -0.0883 \\
\bottomrule
\end{tabular}
\caption*{\footnotesize Identical cell sizes and identical class balance in both conditions; only the direction of time differs.}
\end{table}


\todoitem{Discuss whether the ranking of methods changes between the two
protocols. That comparison, not the leaderboard position, is the main claim of
this report.}

%======================================================================
\section{Discussion}
\label{sec:discussion}

\todoitem{Draw the threads together: what worked, what did not, and what the
two benchmark properties mean for anyone else working on this task.}

\paragraph{Negative results.} \todoitem{List every configuration that did not
clear the three-point threshold, including our own proposals. The
pre-registration commits to reporting these; omitting them is the failure mode
this report is designed to avoid.}

%======================================================================
\section{Threats to validity}
\label{sec:threats}

\paragraph{Internal.} The class--time confound of Section~\ref{sec:confound}
may partly reflect genuine drift in the projects' issue populations rather than
an artefact of sampling. Our claim is restricted to the evaluation protocol:
a random split over temporally structured data yields an optimistic estimate,
regardless of why the structure is there. Separately, 61 test issues (4.1\%)
have a training issue at cosine similarity $\geq 0.90$, and three are exact
duplicates; we report scores with and without that subset.

\paragraph{External.} Five large, well-known projects, with labels balanced
exactly --- a distribution no real issue tracker exhibits, where bug reports
typically dominate. Results should not be read as estimates of deployment
performance on an arbitrary repository.

\paragraph{Construct.} \texttt{question} describes the speech act of the issue
author, whereas \texttt{bug} and \texttt{feature} describe technical content.
The three classes do not lie on a single axis, which plausibly explains why
\texttt{question} is the weakest class in almost every project.

\paragraph{Conclusion.} Small differences on this benchmark are not
measurable; Section~\ref{sec:power} quantifies the threshold, and every claim
in Section~\ref{sec:results} is checked against it.

%======================================================================
\section{Conclusion}
\label{sec:conclusion}

\todoitem{Half a page. Resist the temptation to overstate; the strongest
sentence available is that the benchmark's measurement limits were quantified
and respected.}

%======================================================================
\section*{Contributions}

\todoitem{The lecturer inspects the commit history to see who contributed
what. Fill in this table honestly and make sure it matches
\texttt{git shortlog -sn}.}

\begin{table}[h]
\centering
\begin{tabular}{ll}
\toprule
Member & Contribution \\
\midrule
\todoitem{name} & Track A --- data loading, persistence, preprocessing, EDA \\
\todoitem{name} & Track B --- classical models \\
\todoitem{name} & Track C --- neural models \\
Minh Khanh & Track D --- evaluation, significance, power, benchmark audit \\
\bottomrule
\end{tabular}
\end{table}

%======================================================================
\appendix
\section{Per-seed scores}
\label{sec:appendix-seeds}

\begin{table}[htbp]
\centering
\caption{Raw cross-repository F1 for every seed.}
\label{tab:seeds}
\begin{tabular}{lrrrrrr}
\toprule
Model & 41 & 42 & 43 & 44 & 45 & spread \\
\midrule
linear\_svm & 0.7504 & 0.7504 & 0.7504 & 0.7504 & 0.7504 & 0.00 \\
logreg & 0.7502 & 0.7502 & 0.7502 & 0.7502 & 0.7502 & 0.00 \\
ffnn & 0.7359 & 0.7409 & 0.7314 & 0.7326 & 0.7401 & 0.95 \\
text\_cnn & 0.7191 & 0.7258 & 0.7323 & 0.7319 & 0.7184 & 1.39 \\
random\_forest & 0.7234 & 0.7239 & 0.7242 & 0.7257 & 0.7232 & 0.25 \\
naive\_bayes & 0.7199 & 0.7199 & 0.7199 & 0.7199 & 0.7199 & 0.00 \\
\bottomrule
\end{tabular}
\caption*{\footnotesize Published so that the mean can be checked and so that no seed was quietly dropped. Spread is in F1 points.}
\end{table}


\section{Reproduction}
\label{sec:appendix-repro}

All results are produced by scripts in \texttt{experiments/}, in order:
\texttt{benchmark\_audit.py}, \texttt{02\_baselines.py},
\texttt{03\_setfit.py}, then \texttt{04\_make\_tables.py} to regenerate every
table in this document from \texttt{results/*.json}. Seeds, package versions
and the tuning budget of each model are recorded in those result files.

\begin{thebibliography}{9}

\bibitem{kallis2019tickettagger}
R. Kallis, A. Di Sorbo, G. Canfora, S. Panichella.
\newblock Ticket Tagger: Machine Learning Driven Issue Classification.
\newblock \emph{ICSME}, 2019.

\bibitem{colavito2023fewshot}
G. Colavito, F. Lanubile, N. Novielli.
\newblock Few-Shot Learning for Issue Report Classification.
\newblock \emph{NLBSE}, 2023.

\bibitem{herzig2013misclassification}
K. Herzig, S. Just, A. Zeller.
\newblock It's Not a Bug, It's a Feature: How Misclassification Impacts Bug
  Prediction.
\newblock \emph{ICSE}, pages 392--401, 2013.

\bibitem{benavoli2017bayesian}
A. Benavoli, G. Corani, J. Dem\v{s}ar, M. Zaffalon.
\newblock Time for a Change: a Tutorial for Comparing Multiple Classifiers
  Through Bayesian Analysis.
\newblock \emph{JMLR}, 18(77), 2017.

\todoitem{Add: the NLBSE'24 competition report, Kim (2014) for the CNN, and
the temporal-validation papers named in Section~\ref{sec:related}.}

\end{thebibliography}

\end{document}
